\documentclass[11pt]{article}
\usepackage{amsmath, amssymb}
\usepackage[margin=1in]{geometry}

\newcommand{\dlnr}{\frac{d\ln}{d\ln R}}

\title{Analytical Logarithmic Gradients for the Sirko--Goodman Self-Regulating Disk}
\date{}

\begin{document}
\maketitle

\section{Setup}

In the self-regulating (Sirko \& Goodman) zone, the Toomre parameter is fixed at $Q = Q_{\min}$. The governing equations are (cf.\ \texttt{equations.tex}):
\begin{align}
    \rho &= \frac{\Omega^2}{2\pi G Q_{\min}} \label{eq:rho} \\[6pt]
    c_s &= \left(\frac{\dot{M} f \Omega^2}{6\pi\alpha\rho}\right)^{1/3} \label{eq:cs} \\[6pt]
    P &= c_s^2\,\rho \label{eq:P} \\[6pt]
    H &= \frac{c_s}{\Omega} \label{eq:H} \\[6pt]
    \nu &= \alpha\, c_s\, H \label{eq:visc} \\[6pt]
    \Sigma &= \frac{\dot{M} f}{3\pi\nu} = \frac{\dot{M} f\, \Omega}{3\pi\alpha\, c_s^2} \label{eq:Sigma}
\end{align}
where $\Omega = \sqrt{GM/R^3}$ and $f = 1 - \sqrt{R_*/R}$.

The central temperature $T_c$ is determined by the energy balance (the \texttt{\_fq} equation):
\begin{equation}
    \mathcal{C}_1\, T_c^4 + \mathcal{C}_2\, T_c = \mathcal{C}_3 \label{eq:fq}
\end{equation}
with coefficients:
\begin{align}
    \mathcal{C}_1 &= \frac{4\sigma}{3c} \qquad\text{(constant --- no $R$ dependence)} \label{eq:C1} \\[4pt]
    \mathcal{C}_2 &= \frac{\rho\, k_B}{\mu\, m_p} \label{eq:C2} \\[4pt]
    \mathcal{C}_3 &= \rho\left(\frac{\dot{M} f \Omega^2}{6\pi\alpha\rho}\right)^{2/3}
    = \rho^{1/3}\left(\frac{\dot{M} f \Omega^2}{6\pi\alpha}\right)^{2/3} \label{eq:C3}
\end{align}

\medskip\noindent
\textbf{Key observation:} In this regime, $\rho$, $c_s$, $P$, $H$, $\nu$, and $\Sigma$ are all \emph{explicit} algebraic functions of $R$ (and global constants $M, \dot{M}, \alpha$, etc.). None of them depend on $T_c$. Only $T_c$ itself requires solving the implicit equation~\eqref{eq:fq}.

We define the stored gradients as:
\[
    \texttt{grad\_X} \equiv -\frac{d\ln X}{d\ln R}
\]

%=============================================
\section{Preliminary: Logarithmic Derivatives of Auxiliary Quantities}

\begin{align}
    \frac{d\ln\Omega}{d\ln R} &= -\frac{3}{2} \label{eq:dlnOmega} \\[6pt]
    \frac{d\ln\rho}{d\ln R} &= 2\,\frac{d\ln\Omega}{d\ln R} = -3 \label{eq:dlnrho}
\end{align}

For $f = 1 - \xi$ where $\xi \equiv \sqrt{R_*/R}$:
\begin{equation}
    \frac{df}{dR} = \frac{1}{2}\sqrt{\frac{R_*}{R^3}} = \frac{\xi}{2R}
\end{equation}
\begin{equation}
    \frac{d\ln f}{d\ln R} = \frac{R}{f}\cdot\frac{df}{dR} = \frac{\xi}{2f} \label{eq:dlnf}
\end{equation}

%=============================================
\section{Gradient of $c_s$}

From equation~\eqref{eq:cs}, $c_s^3 = \dot{M} f \Omega^2 / (6\pi\alpha\rho)$, so:
\begin{align}
    3\,\frac{d\ln c_s}{d\ln R}
    &= \frac{d\ln f}{d\ln R} + 2\,\frac{d\ln\Omega}{d\ln R} - \frac{d\ln\rho}{d\ln R} \nonumber\\
    &= \frac{\xi}{2f} + 2\!\left(-\frac{3}{2}\right) - (-3) \nonumber\\
    &= \frac{\xi}{2f} - 3 + 3 = \frac{\xi}{2f}
\end{align}

\begin{equation}
    \boxed{\frac{d\ln c_s}{d\ln R} = \frac{\xi}{6f}}
    \label{eq:dlncs}
\end{equation}

%=============================================
\section{Gradient of $P$ (Pressure)}

From $P = c_s^2\,\rho$:
\begin{align}
    \frac{d\ln P}{d\ln R}
    &= 2\,\frac{d\ln c_s}{d\ln R} + \frac{d\ln\rho}{d\ln R} \nonumber\\
    &= \frac{2\xi}{6f} + (-3) = \frac{\xi}{3f} - 3
\end{align}

\begin{equation}
    \boxed{\texttt{grad\_P} = -\frac{d\ln P}{d\ln R} = 3 - \frac{\xi}{3f}}
    \label{eq:gradP}
\end{equation}

\noindent\textbf{Check:} For $R \gg R_*$, $\xi \to 0$, $f \to 1$, so $\texttt{grad\_P} \to 3$.

%=============================================
\section{Gradient of $\Sigma$ (Surface Density)}

From equation~\eqref{eq:Sigma}, $\Sigma = \dot{M} f\, \Omega / (3\pi\alpha\, c_s^2)$:
\begin{align}
    \frac{d\ln\Sigma}{d\ln R}
    &= \frac{d\ln f}{d\ln R} + \frac{d\ln\Omega}{d\ln R} - 2\,\frac{d\ln c_s}{d\ln R} \nonumber\\
    &= \frac{\xi}{2f} + \left(-\frac{3}{2}\right) - \frac{2\xi}{6f} \nonumber\\
    &= \frac{\xi}{2f} - \frac{\xi}{3f} - \frac{3}{2} \nonumber\\
    &= \frac{3\xi - 2\xi}{6f} - \frac{3}{2} = \frac{\xi}{6f} - \frac{3}{2}
\end{align}

\begin{equation}
    \boxed{\texttt{grad\_Sigma} = -\frac{d\ln\Sigma}{d\ln R} = \frac{3}{2} - \frac{\xi}{6f}}
    \label{eq:gradSigma}
\end{equation}

\noindent\textbf{Check:} For $R \gg R_*$, $\texttt{grad\_Sigma} \to 3/2$.

%=============================================
\section{Gradient of $T_c$ (Central Temperature)}

This is the only quantity that requires implicit differentiation. Start from equation~\eqref{eq:fq}:
\[
    \mathcal{C}_1\, T_c^4 + \mathcal{C}_2\, T_c = \mathcal{C}_3
\]

First, compute the logarithmic derivatives of the coefficients.

\paragraph{$\mathcal{C}_1$:} Constant $\Rightarrow$ $\dfrac{d\ln\mathcal{C}_1}{d\ln R} = 0$.

\paragraph{$\mathcal{C}_2$:} From~\eqref{eq:C2}, $\mathcal{C}_2 = \rho\, k_B/(\mu m_p) \propto \rho$, so:
\begin{equation}
    \frac{d\ln\mathcal{C}_2}{d\ln R} = \frac{d\ln\rho}{d\ln R} = -3
    \label{eq:dlnC2}
\end{equation}

\paragraph{$\mathcal{C}_3$:} From~\eqref{eq:C3}, $\mathcal{C}_3 = \rho^{1/3}\bigl(\dot{M} f \Omega^2/(6\pi\alpha)\bigr)^{2/3}$, so:
\begin{align}
    \frac{d\ln\mathcal{C}_3}{d\ln R}
    &= \frac{1}{3}\,\frac{d\ln\rho}{d\ln R} + \frac{2}{3}\!\left[\frac{d\ln f}{d\ln R} + 2\,\frac{d\ln\Omega}{d\ln R}\right] \nonumber\\
    &= \frac{1}{3}(-3) + \frac{2}{3}\!\left[\frac{\xi}{2f} + 2\!\left(-\frac{3}{2}\right)\right] \nonumber\\
    &= -1 + \frac{2}{3}\!\left[\frac{\xi}{2f} - 3\right] \nonumber\\
    &= -1 + \frac{\xi}{3f} - 2 = -3 + \frac{\xi}{3f}
    \label{eq:dlnC3}
\end{align}

\paragraph{Implicit differentiation.} Taking $d/d(\ln R)$ of both sides of \eqref{eq:fq}:
\begin{equation}
    \underbrace{4\,\mathcal{C}_1\, T_c^4}_{\text{from } \mathcal{C}_1 T_c^4}\cdot g
    + \underbrace{\mathcal{C}_2\, T_c\cdot(-3)}_{\text{from } \frac{d\mathcal{C}_2}{d\ln R}\cdot T_c}
    + \underbrace{\mathcal{C}_2\, T_c}_{\text{from } \mathcal{C}_2\cdot\frac{dT_c}{d\ln R}}\cdot g
    = \mathcal{C}_3\!\left(-3 + \frac{\xi}{3f}\right)
    \label{eq:implicit_full}
\end{equation}
where $g \equiv \dfrac{d\ln T_c}{d\ln R}$.

Solving for $g$:
\begin{equation}
    g\,\bigl(4\,\mathcal{C}_1\, T_c^4 + \mathcal{C}_2\, T_c\bigr)
    = \mathcal{C}_3\!\left(-3 + \frac{\xi}{3f}\right) + 3\,\mathcal{C}_2\, T_c
\end{equation}

\begin{equation}
    \boxed{
        \texttt{grad\_T} = -g = -\frac{\mathcal{C}_3\!\left(-3 + \dfrac{\xi}{3f}\right) + 3\,\mathcal{C}_2\, T_c}{4\,\mathcal{C}_1\, T_c^4 + \mathcal{C}_2\, T_c}
    }
    \label{eq:gradT}
\end{equation}

\paragraph{Limiting cases:}
\begin{itemize}
    \item \textbf{Radiation-dominated} ($\mathcal{C}_1 T_c^4 \gg \mathcal{C}_2 T_c$, i.e.\ $\mathcal{C}_3 \approx \mathcal{C}_1 T_c^4$), with $\xi \to 0$:
    \[
        g \approx \frac{\mathcal{C}_1 T_c^4\cdot(-3) + 0}{4\,\mathcal{C}_1 T_c^4} = -\frac{3}{4}
        \quad\Rightarrow\quad \texttt{grad\_T} \to \frac{3}{4}
    \]
    \item \textbf{Gas-dominated} ($\mathcal{C}_2 T_c \gg \mathcal{C}_1 T_c^4$, i.e.\ $\mathcal{C}_3 \approx \mathcal{C}_2 T_c$), with $\xi \to 0$:
    \[
        g \approx \frac{\mathcal{C}_2 T_c\cdot(-3) + 3\,\mathcal{C}_2 T_c}{\mathcal{C}_2 T_c} = 0
        \quad\Rightarrow\quad \texttt{grad\_T} \to 0
    \]
\end{itemize}

%=============================================
\section{Summary}

All gradients for the self-regulating zone, with $\xi = \sqrt{R_*/R}$ and $f = 1 - \xi$:

\begin{center}
\renewcommand{\arraystretch}{2.0}
\begin{tabular}{l c c}
\hline
\textbf{Gradient} & \textbf{Formula} & \textbf{Limit ($R \gg R_*$)} \\
\hline
$\texttt{grad\_T} = -\dfrac{d\ln T_c}{d\ln R}$
    & Eq.~\eqref{eq:gradT} (implicit) & $3/4$ \\[6pt]
$\texttt{grad\_Sigma} = -\dfrac{d\ln\Sigma}{d\ln R}$
    & $\dfrac{3}{2} - \dfrac{\xi}{6f}$ & $3/2$ \\[6pt]
$\texttt{grad\_P} = -\dfrac{d\ln P}{d\ln R}$
    & $3 - \dfrac{\xi}{3f}$ & $3$ \\
\hline
\end{tabular}
\end{center}

Note that $\texttt{grad\_P} = 2\,\texttt{grad\_Sigma}$ for all $R$, which follows from $P = c_s^2\rho$ and $\Sigma \propto f\Omega/c_s^2$, giving $P/\Sigma \propto c_s^4 \rho / (f\Omega) \propto \Sigma$ after simplification. This provides a useful cross-check.

\end{document}